\subsection{Una casilla marcada y el incremento del determinante [D]}
Queda el incremento de $R_{\rm det}$. Se puede medir el efecto de $e=(r,c)$ sobre $L_T$
variando su peso o comparando el tablero con el menor tras ocuparla. Para $e$ permitida, sean
$\pi^{\rm der}_e=\partial_{\log w_e}L_T=p_e(1+\hat c_e)$,
$\varepsilon_e=p_e/P_e-1$ y
$r_e=L_T(M_e)-L_T(M)-\log(\pi^{\rm der}_e/w_e)$.
El Teorema A de torres da $|\hat c_e|\le C/k$ y $|r_e|\le C/k^2$. Así $r_e$ mide hasta
segundo orden la discrepancia entre ambas operaciones. Usando el escalado
del menor $P'=\operatorname{diag}(x')P^{(e)}\operatorname{diag}(y')$, donde $P^{(e)}$ es la
submatriz de $P$, se obtiene la identidad exacta
\begin{align*}
 R_{\rm det}(M_e)-R_{\rm det}(M)&=-r_e-\mathsf E_e,\\
 \mathsf E_e={}&\log(1+\varepsilon_e)+\log(1+\hat c_e)
          +\sum_{i\ne r}\log x'_i+\sum_{l\ne c}\log y'_l\\
 &+(k-1)\log(1-1/k)
  +\tfrac12\log\frac{\det(I-b_{P'}^\top b_{P'})}{\det(I-b_P^\top b_P)}\\
 &+\frac{\delta_1((k-1)P')}{k-1}-\frac{\delta_1(kP)}k .
\end{align*}
En efecto, $\Delta=-L_T(P)+\log(k!/k^k)-\frac12\log\det(I-b_P^\top b_P)$.
La covariancia de $L_T$ bajo escalado da las dos sumas de potenciales. El signo del término de
tamaño sale del cociente de los dos factoriales. Para demostrar $(A^+)$ debemos ganar una potencia
de $k$ en la suma firmada $r_e+\mathsf E_e$. La expresión exacta de $L_T$ conserva el
desplazamiento cuadrático del punto estacionario, habitualmente denotado por $\zeta$. Acotar los
once restos de segundo orden en valor absoluto perdería su cancelación.

\paragraph{Parámetros independientes.}
Para conservar la suma firmada, situamos los dos tableros en una misma familia. Variamos por
separado la masa de la fila $r$, la de la columna $c$, el tamaño inverso y el peso de la casilla
$e$. Fije $W=kP$ y tome cada fila y columna como un tipo. Con $u=1/k$, sean $A_-,B_-$ las
probabilidades uniformes que omiten respectivamente la fila y la columna distinguidas. Ponga
\[
  A_s=(1-s)A_-+s\delta_r,\qquad B_t=(1-t)B_-+t\delta_c,\qquad h=1/N .
\]
Multiplique sólo $W_{rc}$ por $v$. En estas medidas los núcleos escalado y estacionario resuelven
$N=W^v\exp(a_i+b_l)$ y $q(1-hq)=W^v\exp(z_i+y_l)$ con marginales uno.
La inversa de la derivada, fijado el calibre, está uniformemente acotada en norma supremo.
La norma cuadrática se controla por la brecha, y la parte integral transforma norma cuadrática
en supremo por $W\le M_0$. Perturbar las medidas cuesta $O(|s-u|+|t-u|)$, y el peso
marcado cuesta $O(\max(|s|,|t|))$ para $|v|<2$. La iteración con la inversa del centro
converge en discos fijos de clase para $h,s,t$ y da la rama holomorfa.
Los determinantes de los núcleos centrados $V$ que intervienen admiten la forma
\[
 D(V)=\sum_{\ell\ge0}\frac{(-1)^\ell}{(\ell!)^2}
  \int\det[V(i_a,l_b)]_{a,b\le\ell}^{\,2}\,
               dA_s^{\otimes\ell}dB_t^{\otimes\ell}.
\]
Cauchy--Binet y Hadamard dan un mayorante $\ell^\ell C^\ell/(\ell!)^2$. La brecha y la
norma de Hilbert--Schmidt del centro dan un suelo positivo para el determinante allí.
La serie y Cauchy permiten tomar su logaritmo en discos uniformemente menores. En la expresión
de $\Delta$ se sustituye $\log(N!/N^N)-\frac12\log(4\pi N)$ por
$-1/h-\frac12\log2+h/12$ y se guarda aparte el resto $\eta_N=O(N^{-3})$ de Stirling.
La expresión resultante $D_*$ tiene $h^2D_*$ acotada y holomorfa en esos discos.
Para masas racionales positivas, replicar cada tipo $N A_s(i)$ o $N B_t(l)$ veces
preserva el operador de tipos y añade sólo valores singulares cero. Aplicando la parte
determinista de G1 a esas réplicas queda $D_*=\delta_1/N+O(N^{-2})$. Así los tres primeros
coeficientes de $h^2D_*$ se anulan y el de grado tres es $\delta_1$. Por continuidad e identidad
holomorfa, estos coeficientes coinciden para todas las masas. Si $|h|\le r/2$ en un disco
de radio fijo $r$ donde $|h^2D_*|\le M$, el resto tras el grado tres, dividido por $h^4$,
está acotado por $\sum_{j\ge4}Mr^{-j}(r/2)^{j-4}=2M/r^4$. Resulta
$R_{\rm det}=\eta_N+h^2D_2(h,s,t,v)$, con las derivadas primeras de $D_2$ acotadas por
la misma fórmula integral en discos interiores.
Sobre una cara de masa cero se replica sólo la parte visible.
Su calibre se elimina en las sumas de potenciales y se identifica su propio $\delta_1$.

El paso al menor sustituye $(h,s,t)=(u,u,u)$ por $(u/(1-u),0,0)$. La variación de $h$
es $O(u^2)$ y la de cada masa es $O(u)$, de donde $\Delta_e R_{\rm det}=O(u^3)$.
La identidad de arriba prueba $(A^+)$ con el signo indicado. En la cara $s=0$ o $t=0$ las
expresiones integradas son independientes de $v$. Por ello el desarrollo de $\partial_vD_2$
carece de los términos con potencia cero de $s$ o $t$. Dividir su serie por $st$ da una
función $H$ acotada en discos interiores por la serie geométrica de las cotas de Cauchy.
En $s=t=u$ la derivada de $R_{\rm det}$ respecto a $v$ es $O(u^4)$ hasta $v=0$.
Integrar sobre $[0,1]$ da el borrado de $(D^+)$. Al sumar los incrementos de $(L^{*+})$ se
obtiene $(T^+)$ para el resto total $R$. Así puede compararse la aproximación de un mismo par
antes y después de pasar a un menor adicional que lo conserve.
